% Шаблон тезисов доклада на научном семинаре
% Лаборатория математики ФН1
% Компиляция: pdflatex abstract.tex && pdflatex abstract.tex
% Объём — не более одной страницы

\documentclass[12pt, a4paper]{article}

\usepackage[T2A]{fontenc}
\usepackage[utf8]{inputenc}
\usepackage[russian]{babel}
\usepackage[margin=20mm]{geometry}
\usepackage{amsmath, amssymb, amsthm}

\pagestyle{empty}
\setlength{\parindent}{1.25cm}

\theoremstyle{plain}
\newtheorem{theorem}{Теорема}

\begin{document}

\noindent УДК 517.9

\begin{center}
  {\large\bfseries Название доклада}

  \vspace{0.4em}

  \textit{Фамилия И.\,О.}

  \vspace{0.2em}

  {\small Лаборатория математики ФН1\\
  \texttt{email@example.ru}}
\end{center}

\vspace{0.8em}

\noindent\textbf{Аннотация.} Пять--семь предложений: какая задача рассматривается,
каким методом она решается и что получено. Аннотация должна быть понятна
слушателю, не работающему в этой области.

\vspace{0.5em}

\noindent\textbf{Ключевые слова:} три--пять терминов через запятую.

\vspace{1em}

Рассматривается задача
\begin{equation}\label{eq:problem}
  L u = f, \qquad u \in D(L),
\end{equation}
где $L$~--- линейный оператор, действующий в гильбертовом пространстве $H$.

Известно, что при выполнении условия $\ldots$ задача~\eqref{eq:problem}
однозначно разрешима. В настоящей работе это условие ослаблено.

\begin{theorem}
  Формулировка основного результата.
\end{theorem}

Доказательство опирается на $\ldots$ Приводится численный эксперимент,
подтверждающий полученную оценку.

\vspace{1em}

\begin{thebibliography}{9}
  \footnotesize
  \bibitem{ref1} \emph{Автор А.\,Б.} Название статьи~// Название журнала.~---
    2024.~--- Т.~10, №~2.~--- С.~15--28.
  \bibitem{ref2} \emph{Автор В.\,Г.} Название книги.~--- М.: Издательство, 2023.~--- 320~с.
\end{thebibliography}

\end{document}
